\section{Where the single schema cracks}
\label{sec:ch01-cracks}

What follows is not a list of things that might go wrong. Each item has a
company's name and a date on it, because \enquote{the monolith was fine} is an
argument I have had several times, and it is easier to have with sources.

\index{schema!review contention}%
Start with the review queue, because it is the first thing anybody notices.
Booking.com described their pre-federation GraphQL service plainly: as it grew,
one team was left with \enquote{dozens of merge requests a day to review for
schema changes and service mappings}, while also managing frequent deployments
of that service and helping resolve merge conflicts between
teams~\autocite{ernst2022booking}. The detail I find most useful is what this did
to ownership. In their words, it \enquote{made it difficult for service owners
to have strong ownership of their data and schema as they were heavily tied to
the single service}. The teams that understood a domain were not the teams that
could change how it was exposed. Ownership had become a label rather than a
capability.

\index{API monolith}%
The second failure mode is the one I would put on a poster, and Netflix has the
clearest telling of it because they walked into it having already done the hard
work of avoiding it. They had broken their system into microservices, and then,
as Jennifer Shin put it in a 2020 talk, found that those services were
\enquote{all aggregated into a single API monolith}. Her colleague Stephen
Spalding described how it happened: the gateway that composed all those services
kept acquiring responsibilities, fallback logic to handle failure gracefully,
then caches, then in-memory datastores, then business logic, until \enquote{the
API gateway had become the new monolith}~\autocite{shin2020netflixtalk}.

That is worth sitting with, because nothing in it is about GraphQL. Any
aggregation layer with one owner and many contributors drifts toward becoming
the system's centre of gravity. GraphQL just makes the drift legible, since the
accumulated coupling has a name and a file you can open.

\index{schema!quality decay}%
Third, schema quality decays when the people who own the schema are not the
people who own the domain. Netflix's account of their Studio API bottlenecks
names three causes: the API team was \enquote{disconnected from the domain
expertise and the product needs, which negatively impacted the schema's
health}; wiring a new back-end element into the graph was manual work that
\enquote{ran counter to the rapid evolution promised by a microservice
architecture}; and one small team could not keep up with the operational and
support burden of a growing graph~\autocite{shikhare2020netflix1}. The
consequence they report is the kind of thing that never appears in an
architecture diagram: inconsistent data across Studio applications was the top
support issue in their studio engineering organisation in 2018. A modelling
problem had turned into a support queue, which is usually how modelling problems
announce themselves.

\index{deploy coupling}%
Fourth, deploy coupling, which in practice means two teams have to ship before
one field exists. Tejas Shikhare put the fix in terms of the problem: with
federation, \enquote{you don't have to implement the feature in the backend
service, in the GraphQL service before the client can use it}, so the team that
owns the data can ship it and it is simply
there~\autocite{shikhare2023podcast}. Volvo Car Mobility described the same
tax from the other side: every schema change meant working inside the monolith
\emph{and} defining service-to-service APIs \emph{and} changing one or more
microservices, and on top of that their uptime depended on that single
service~\autocite{radjavi2023volvo}.

\index{breaking changes}%
Fifth, and this one is quieter, you eventually cannot change anything because
you do not know who is reading it. Matt Oliver of Major League Baseball
described having \enquote{low visibility into who is making calls on our
platform}, with clients owned by different teams, which made model changes and
breaking changes hard, because notifying and coordinating those clients was
itself the difficult part~\autocite{oliver2022mlb}. I want to flag an honest
caveat here: federation does not fix this. What fixes it is a schema registry
and operation-level tracking, which federated tooling tends to bring along
because composition needs a registry anyway. You can have that without
federating, and chapter~\ref{ch:do-you-actually-need-federation} says so.

\index{velocity ceiling}%
The cumulative version is a ceiling that hiring does not raise. Wayfair's Leah
Hurwich Adler, describing their monolith to Apollo, said it was \enquote{so
unwieldy that we hit a ceiling in our velocity, and hiring more engineers didn't
help us push code to production faster}~\autocite{adler2024wayfair}. That
sentence is the shape of the whole problem. If throughput were limited by
engineering capacity, more engineers would help. It was limited by a shared
resource, and adding contenders for a shared resource makes the contention
worse.

The most candid summary I know comes from PayPal in 2019, early enough that the
author was not describing a solved problem. Mark Stuart wrote that users
reasonably want \enquote{a single, cohesive graph that they can traverse through
without thinking about many services}, and then immediately conceded: \enquote{In
reality, assembling a single graph is difficult.} He calls it \enquote{the most
difficult problem that we have with GraphQL}, and locates the difficulty
correctly: \enquote{scaling in the Enterprise isn't about horizontal scaling or
paying a lot of money for servers or cloud compute. Scaling people, tooling and
processes are most challenging}~\autocite{stuart2019paypal}.

Now read back over that list and notice what is missing from it. Not one of
these is a performance problem, and not one is a limitation of GraphQL as a
language. They are all questions about who is allowed to change what, and how
long everyone else waits while they do. That is a specific kind of problem with
a specific and well known cause, which is the next section.
